\abstract{РЕФЕРАТ}

Объем работы равен \formbytotal{lastpage}{страниц}{е}{ам}{ам}. Работа содержит \formbytotal{figurecnt}{иллюстраци}{ю}{и}{й}, \formbytotal{tablecnt}{таблиц}{у}{ы}{}, \arabic{bibcount} библиографических источников и \formbytotal{числоПлакатов}{лист}{}{а}{ов} графического материала. Количество приложений – 2. Графический материал представлен в приложении А. Фрагменты исходного кода представлены в приложении Б.

Перечень ключевых слов: информационная система, автоматизация, центр повышения квалификации, база данных, веб-приложение, Flask, PostgreSQL, Python, Tkinter, импорт данных, фильтрация, слушатель, отчётность.

Объектом разработки является программно-информационная система для автоматизации деятельности учебно-методического центра повышения квалификации университета.

Целью выпускной квалификационной работы является создание программно-информационной системы, обеспечивающей централизованное хранение данных о слушателях, гибкий поиск и фильтрацию, импорт данных из Excel, формирование статистической отчётности и разграничение прав доступа.

В процессе создания системы были проанализированы особенности деятельности центра, выделены основные сущности предметной области, спроектирована и реализована реляционная база данных (PostgreSQL). Разработано веб-приложение на микрофреймворке Flask с использованием шаблонизатора Jinja2 и библиотеки Flask-Login для аутентификации. Для управления серверной частью создано десктоп-приложение на Tkinter, которое запускает и останавливает встроенные серверы PostgreSQL (библиотека pgembed) и Flask.

Система прошла системное тестирование, подтвердившее соответствие всем требованиям технического задания. Разработанная программно-информационная система может быть использована в учебно-методических центрах повышения квалификации для автоматизации учёта слушателей и подготовки отчётности.

\selectlanguage{english}
\abstract{ABSTRACT}
  
The volume of work is \formbytotal{lastpage}{page}{}{s}{s}. The work contains \formbytotal{figurecnt}{illustration}{}{s}{s}, \formbytotal{tablecnt}{table}{}{s}{s}, \arabic{bibcount} bibliographic sources and \formbytotal{числоПлакатов}{sheet}{}{s}{s} of graphic material. The number of applications is 2. The graphic material is presented in annex A. The layout of the site, including the connection of components, is presented in annex B.

List of keywords: information system, automation, professional development center, database, web application, Flask, PostgreSQL, Python, Tkinter, data import, filtering, student, reporting.

The object of development is a software information system for automating the activities of a university professional development center.

The purpose of the final qualifying work is to create a software information system that provides centralized storage of student data, flexible search and filtering, data import from Excel, generation of statistical reports, and access control.

During the system development, the specifics of the center's activities were analyzed, the main entities of the subject area were identified, and a relational database (PostgreSQL) was designed and implemented. A web application was developed using the Flask microframework with the Jinja2 templating engine and the Flask-Login library for authentication. A desktop application based on Tkinter was created to manage the server side, which starts and stops the embedded PostgreSQL server (using the pgembed library) and the Flask server.

The system has passed system testing, confirming compliance with all requirements of the technical specification. The developed software information system can be used in professional development centers to automate student record keeping and report generation.
\selectlanguage{russian}
